%
% 22-fortsetzung.tex -- Analytische Fortsetzung entlang einer Kurve
%
% (c) 2026 Prof Dr Andreas Müller
%
\subsection{Analytische Fortsetzung entlang einer Kurve}
Sei $\gamma$ eine Kurve in der komplexen Ebene und $f$ eine
Funktion, die in einer Umgebung des Punktes $\gamma(0)$ der
Kurve holomorph ist.
Die Funktion ist also in einem Kreis vom Radius $\varrho(\gamma(0))$
um den Punkt $\gamma(0)$ analytisch.
Nach \eqref{buch:analytisch:fortsetzung:eqn:taylor01}
kann die Funktion $f(z)$ in jedem Punkt $\gamma(t)$ wieder in eine
Potenzreihe mit Konvergenzradius $\varrho(\gamma(t))$ entwickelt
werden, so dass die beiden Potenzreihen auf dem Schnittgebiet der
beiden Konvergenzkreise übereinstimmen.
Für einen Punkt $\gamma(t)$ nahe dem Rand des Konvergenzkreises 
der Potenzreihe im Punkt $\gamma(0)$ kann der Konvergenzkreis
der neuen Reihe über den Rand des alten Konvergenzkreises hinausragen.
Die neue Potenzreihe erweitert also die Funktion über den ursprünglichen
Konvergenzkreis hinaus.
Indem dies entlang der Kurve immer wieder wiederholt wird, kann die
Funktion auf ein Umgebung der Kurve ausgedehnt werden, wie dies zum
für die Wurzelfunktion in
Abbildung~\ref{buch:analytisch:fortsetzung:fig:wurzel2}
illustriert wird.

\begin{beispiel}
Die Funktion $f(z)=\frac{1}{z}$ ist holomorph in $\mathbb{C}\setminus\{0\}$.
Sie $\gamma(t)=e^{2\pi it}$ der Einheitskreis.
Im Punkt $\gamma(0)=1$ kann die Funktion durch die geometrische Reihe
\begin{equation}
\frac{1}{z}
=
\frac{1}{1+(z-1)}
=
\sum_{k=0}^\infty (z-1)^k
\label{buch:analytisch:fortsetzung:bsp:geom1}
\end{equation}
dargestellt werden, die in $B_1(1)$ konvergent ist.
Der Punkt $\gamma(\frac{\pi}4)$ liegt innerhalb dieses Konvergenzkreises.
Die Entwicklung
\begin{equation}
\frac{1}{z}
=
\frac{1}{a+(z-a)}
=
\frac1a \frac{1}{1+\displaystyle\frac{z-a\mathstrut}{a}}
=
\sum_{k=0}^\infty
\frac{1}{a^k}(z-a)^k
\label{buch:analytisch:fortsetzung:bsp:geoma}
\end{equation}
ist eine Potenzreihenentwicklung im Punkt $a$ mit Konvergenzradius $|a|$.
Für $a=\gamma(\frac{\pi}4)$ ist die Reihe von
\eqref{buch:analytisch:fortsetzung:bsp:geoma},
die in einer Umgebung von $a$ mit der Reihe
\eqref{buch:analytisch:fortsetzung:bsp:geom1}
übereinstimmt.
Der Punkt $i\in\mathbb{C}$ ist nicht im Konvergenzkreis der Reihe 
\eqref{buch:analytisch:fortsetzung:bsp:geom1},
der Wert $f(i)$ kann also damit nicht berechnet werden,
wohl aber mit
\eqref{buch:analytisch:fortsetzung:bsp:geoma}.
\end{beispiel}

Im Beispiel führt die analytische Fortsetzung entlang der Kurve 
nach einem Umlauf wieder zur selben Entwicklung
\eqref{buch:analytisch:fortsetzung:bsp:geom1}
im Punkt $1$.
Dies ist jedoch nicht selbstverständlich, wie die Beispiele
im kommenden Abschnitt zeigen werden.
